\documentclass[paper=a4, fontsize=11pt]{jlreq} % 日本語用の標準的なクラス

% --- パッケージ群 ---
\usepackage[utf8]{inputenc}
\usepackage[T1]{fontenc}
\usepackage{amsmath, amssymb, amsthm} % 数学の基本
\usepackage{physics} % \ket{0}, \pdv{x} などが簡単に打てる
\usepackage{bm} % 太字ベクトル \bm{v}
\usepackage{graphicx}
\usepackage{hyperref} % リンク・目次用
\usepackage{cite} % 引用用

% --- 定理環境の設定 (OMCの問題風に整理できる) ---
\theoremstyle{definition}
\newtheorem{theorem}{定理}[section]
\newtheorem{definition}[theorem]{定義}
\newtheorem{example}[theorem]{例}
\newtheorem{remark}[theorem]{注釈}
\newtheorem{problem}[theorem]{問題} % 自作問題用

% --- 独自のショートカット (CFT/パンルヴェ用) ---
\newcommand{\Virexp}[2]{[L_{#1}, L_{#2}]} % ヴィラソロの交換関係
\newcommand{\taufunc}{\tau(z)} % タウ関数

\title{共形場理論とパンルヴェ方程式に関する研究ノート}
\author{松浦光佑}
\date{\today}

\begin{document}

\maketitle
\tableofcontents % 自動で目次を作成

\section{はじめに}
本ノートは、岡本『パンルヴェ方程式』に関する学習記録である。
日々の学習を「AC（Accepted）」として記録し、GitHubでの進捗可視化を目的とする。

\section{パンルヴェ方程式の基礎}
ここに、今日思い出した定義を一つだけ書いてみる。
\begin{definition}[リーマン図式]
    % --- ここに数式 ---
    \[
    \left\{
    \begin{matrix}
        x = \xi_1 & x = \xi_2 & \cdots & x = \xi_m & x = \infty \\
        s_1^{(1)} & s_1^{(2)} & \cdots & s_1^{(m)} & s_1^{(\infty)} \\
        s_2^{(1)} & s_2^{(2)} & \cdots & s_2^{(m)} & s_2^{(\infty)} \\
        \vdots    & \vdots    &        & \vdots    & \vdots \\
        s_n^{(1)} & s_n^{(2)} & \cdots & s_n^{(m)} & s_n^{(\infty)}
    \end{matrix}
    \right\}
    \]
    \begin{itemize}
        \item[\textbf{Type:}] 確定特異点$\xi_{i}$と各確定特異点の特性指数$s_j^{(i)}$を書き並べたもの。
        \item[\textbf{Insight:}] この表の数値（局所的な性質）を\textbf{一切変えずに}、特異点の位置 $\xi_i$ だけを時間発展させた時に、裏で満たされるべき非線形の方程式が「ガルニエ系」。つまり、モノドロミー保存変形において\textbf{「絶対に定数として保存したいターゲット」}がこれ。
        \item[\textbf{Link:}] 岡本2.4節の(2.81)/p.81を参照。
    \end{itemize}
\end{definition}



% --- 今後の課題（TODOリスト） ---
\section{Next Step / To-Do}
\begin{itemize}
    \item [ ] ヴィラソロ代数の中心電荷 $c$ の物理的意味の解釈。
    \item [ ] パンルヴェIIのタウ関数を具体的にTeXで書き出す。
\end{itemize}

\end{document}